\section{Introduction}

Enhancing the capabilities of (LLMs) in the realms of mathematical reasoning and autoformalization is essential for advancing the current state of artificial intelligence and automated theorem proving. The goal to enhance LLMs with a generative capacity for mathematical proofs is a necessity for the evolution of computational reasoning. Current LLM models demonstrate an impressive understanding over a large range of language tasks, yet their proficiency in deciphering and formulating the more rigorous and nuanced formal mathematics is not consistently up-to-par. This deficiency underscores a critical gap in the foundational understanding of mathematical logic and its structural nuances within contemporary models, a gap that existing methodologies have scarcely managed to bridge \cite{polu2022formal}.

The challenge at hand encompasses the transformation of LLMs from mere processors of mathematical data to more complete tools capable of engaging with, understanding, and generating formal mathematical proofs. Traditional approaches to this problem have primarily revolved around direct training with mathematical datasets or integrating elementary problem-solving tasks. While these strategies have yielded small improvements in understanding model performance of different sizes, they fall short of addressing the complex semantic and logical frameworks that are core to formal mathematical proofs. Such methodologies, therefore, leave a large gap in the models' capacity for formal reasoning and the autoformalization of mathematical statements, highlighting the limitations of current patterns in achieving such comprehension and generation of formal mathematical content \cite{bansal2019holist}.

In light of these challenges, our research proposes a paradigmatic shift through the direct leverage of verified proof libraries from Interactive Theorem Provers (ITPs) such as Lean, Coq, and Isabelle. We hypothesize that mining step-by-step formal proofs from these libraries and integrating them into the training corpus of LLMs can catalyze a profound enhancement in their ability to understand and generate formal mathematical proofs. This innovative approach is predicated on the potential of ITPs to provide a structured and rigorous foundation for learning, one that is meticulously aligned with the logical precision and semantic depth required for formal mathematics. Through this methodology, we aim to furnish LLMs with a robust framework for engaging with formal proofs, thereby surmounting the prevailing barriers in mathematical reasoning and autoformalization \cite{szegedy2020promising} \cite{wu2022autoformalization}.

To address this gap, we have created a new pipeline strategy centered around the fine-tuning of a Mistral-7B (7 billion) parameter model. The main part of this process involves prompting the model to reinterpret formal mathematical statements into informal statements of ``top level machine learning research quality''. This plan seeks to form a bridge between the rigorous domain of formal proofs and the more intuitive, natural language processing capabilities of LLMs. The resultant dataset from our Mistral pipeline is represented in a structured JSON tree of formal and informal proofs and will be used to fine-tune several of our models (Mistral-7B, Gemma-7B, GPT-3.5) to measure improvements \cite{yang2023leandojo} \cite{jiang2023draft}.

The evaluation of our methodology's effectiveness will be anchored in a comprehensive benchmarking strategy. We will employ the Massive Multitask Language Understanding (MMLU) benchmark to assess enhancements in general intelligence and mathematical reasoning, supplemented by the ProofNet benchmark for evaluating autoformalization skills. We will additionally test on different subsections of the Hendrycks MATH dataset. These benchmarks are specifically designed to measure the nuanced capabilities of LLMs in processing and generating mathematical content, thereby providing a holistic view of the models' proficiency in the post-enhancement landscape. Additionally, as a pre-screen check, we will manually trace through our generated dataset to see if it makes sense to a human first, to increase our confidence of having effective and correct data passed to our Mistral model. Through this multifaceted evaluation framework, we aim to test our hypothesis that formal proofs can significantly elevate LLMs' reasoning and understanding, allowing for new improvements in automated reasoning and theorem proving \cite{hendrycks2021measuring} \cite{polu2020generative}.

Notwithstanding the promise of our approach, we acknowledge the inherent challenges in enabling LLMs to generate a diverse array of informal proofs from formal statements. The complexity and variations of mathematical proofs imply that a singular formal statement can engender multiple informal interpretations, posing a formidable challenge in ensuring the consistency and accuracy of the model's outputs. Addressing this challenge necessitates not only sophisticated model architecture and training methods but also a deeper understanding of the interplay between formal and informal mathematical reasoning.

The implications of successfully navigating these challenges extend beyond the immediate advancements in LLM capabilities. By fine-tuning a model that improves upon understanding of formal mathematics, we hope to contribute to developments in automated theorem proving and artificial general intelligence. This research contains the potential to redefine the application of LLMs in formal mathematics, offering significant implications for a large range of fields reliant on mathematical reasoning and logic. Furthermore, it sets the stage for future explorations into the nuanced relationship between formal proofs and natural language processing, potentially sparking innovations in other areas of AI where formal reasoning plays a critical role \cite{wu2022autoformalization_arxiv} \cite{zheng2022minif2f}.

In short, our research aims to form a stronger connection between LLMs and formal mathematics through methods that leverage the rigor and strictness of Interactive Theorem Provers (ITPs) and existing large language models' NLP capabilities. By enhancing LLMs' capacity to reason mathematically and autoformalize proofs, we aim to make advancements in automated reasoning, artificial general intelligence, and more.
